\subsection{SuperCollider}\label{subsec:supercollider}

\textbf{SuperCollider} is a programming language and environment for real-time audio synthesis and algorithmic
composition, originally developed by James McCartney.
It was designed to provide a highly expressive and flexible platform for both sound synthesis and musical structure
definition, combining a high-level language with a powerful synthesis engine~\cite{McCartney_2002}.

Unlike traditional music programming environments, \textbf{SuperCollider} was built with the goal of supporting dynamic
and generative compositions, where musical processes can evolve over time rather than representing a fixed score.
It enables composers to describe not only concrete musical events, but also systems that generate music
algorithmically~\cite{McCartney_2002}.

%% ---------------------------------------------------------------------------------------------------------------------

\subsubsection{Key Features}

\textbf{SuperCollider} follows an \textit{object-oriented and functional programming model}, where all
elements are treated as objects and can be composed using high-level abstractions.
The language provides a rich set of programming constructs, including functions, control structures, and data types,
allowing for complex algorithmic manipulation of sound~\cite{McCartney_2002}.

A central abstraction in \textbf{SuperCollider} is the concept of a \textit{unit generator} (\textit{UGen}),
which represents a basic signal processing component.
These UGens can be combined to form synthesis graphs, defining how audio signals are generated and processed.
Instruments are constructed functionally, meaning that a program describes how to build a network of UGens rather than
specifying a fixed structure~\cite{McCartney_2002}.

Example snippet:

\begin{lstlisting}[label={lst:supercollider-1}]
  {
    SinOsc.ar(440, 0, 0.2)
  }.play;
\end{lstlisting}

\textbf{SuperCollider} also introduces higher-level abstractions such as \textit{patterns} and \textit{streams}, which
allow composers to define sequences of musical events.
A stream represents a sequence of values generated over time, while patterns define reusable structures that can produce
multiple streams~\cite{McCartney_2002}.

Another important architectural feature is the separation between the \textit{language} and the
\textit{synthesis server}.
The language is responsible for generating musical events, while the synthesis engine executes them.
Communication between the two is handled through the \textit{Open Sound Control (OSC)} protocol, enabling distributed
and real-time control over audio processes~\cite{McCartney_2002}.

This architecture allows \textbf{SuperCollider} to support:

\begin{itemize}
  \item real-time audio synthesis
  \item dynamic creation and modification of sound processes
  \item distributed and network-based performance setups
\end{itemize}

%% ---------------------------------------------------------------------------------------------------------------------

\subsubsection{Limitations}

Despite its expressive power, \textbf{SuperCollider} presents several challenges.

First, the language has a \textit{steep learning curve}, largely due to its combination of object-oriented, functional,
and domain-specific paradigms.
Understanding concepts such as UGens, synthesis graphs, and event streams requires both programming knowledge and
familiarity with digital signal processing.

Second, the separation between the language and the synthesis server introduces additional complexity.
Communication via \textit{OSC} can lead to latency and requires users to reason about asynchronous execution
and message passing~\cite{McCartney_2002}.

Furthermore, while \textbf{SuperCollider} excels at sound synthesis and real-time performance, it is less focused on
structured, high-level composition.
Musical ideas are often expressed through low-level signal processing constructs, which can make it harder to represent
abstract musical structures such as timelines or hierarchical compositions.

%% ---------------------------------------------------------------------------------------------------------------------

\subsubsection{Relevance to Our DSL}

\textbf{SuperCollider} demonstrates a powerful approach to integrating sound synthesis with a high-level
programming language, offering extensive flexibility in defining and manipulating audio processes.
Its use of composable abstractions, such as unit generators and event streams, provides valuable insight into designing
expressive musical systems.

However, its focus on low-level synthesis and real-time interaction contrasts with the goals of this thesis,
which emphasize structured composition and deterministic behavior.
The proposed DSL aims to provide a higher-level representation of musical concepts, enabling clearer expression of
temporal relationships and compositional structure, while still allowing integration with synthesis backends.
